\chapter{PEP Post Exposure Prophylaxis}
\index{PEP Post Exposure Prophylaxis}

\medskip
\noindent\textit{\textbf{Under review.} This chapter is an AI-assisted educational draft; it is not a substitute for professional medical advice.}

\section*{Definition and Overview}
a short course of antiretroviral medicines started after a possible HIV exposure to reduce the chance of infection.

\section*{Clinical Features}
there are usually no immediate symptoms from the exposure or from early HIV infection; anxiety and medicine side effects such as nausea may occur.

\section*{When to Seek Medical Care}
Seek medical review for new, persistent, worsening, or function-limiting symptoms. Contact your local emergency services for severe breathing difficulty, collapse, sudden neurological change, severe chest or abdominal pain, or any rapidly worsening illness.

\section*{Diagnosis and Investigations}
urgent risk assessment considers exposure type, timing, source status, pregnancy, kidney function, and baseline HIV and STI tests.

\section*{Management}
start as soon as possible, ideally within hours and generally no later than 72 hours, and complete the prescribed course with follow-up testing. Seek urgent sexual-health or emergency care.

\section*{Complications}
Complications depend on the underlying cause and may include progression, effects on other organs, treatment adverse effects, or reduced quality of life. Early assessment can reduce preventable harm.

\section*{Prognosis and Self-Care}
Outcome depends on the cause, severity, complications, and response to treatment. Follow the care plan, keep appointments, avoid known triggers, and seek help if symptoms worsen.
\section*{Safety-netting}
Seek urgent help for severe or rapidly worsening symptoms, collapse, breathing difficulty, severe pain, confusion, dehydration, uncontrolled bleeding, or any immediate danger.
